\documentclass[11pt]{article}
\usepackage[final]{acl}
\usepackage{times}
\usepackage{latexsym}
\usepackage{fontspec}
\usepackage{polyglossia}
\setdefaultlanguage{russian}
\setotherlanguage{english}
\setmainfont{Times New Roman}
\newfontfamily\cyrillicfont{Times New Roman}
\usepackage{microtype}
\usepackage{inconsolata}
\usepackage{graphicx}
\usepackage{booktabs}

\title{Question Answering}

\author{Игорь Пивнев, ИТМО}

\begin{document}
\maketitle

\begin{abstract}
В работе исследуются методы автоматического ответа на вопросы на коллекции MIRAGE.
Рассматриваются два подхода: экстрактивный QA на основе модели RoBERTa и
генеративный QA на основе инструкционной языковой модели Qwen2.5-1.5B-Instruct.
Эксперименты проводятся в режимах oracle (верный документ), top-1 и top-5
(конкатенация фрагментов) с использованием двух ранкеров: BM25 и CrossEncoder.
Для оценки используются метрики EM, F1 (SQuAD-стиль), EM-strict, EM-loose
(MIRAGE-стиль) и BERTScore. Дополнительно применяется подход LLM-as-a-judge
для анализа качества ответов.
\end{abstract}

\section{Постановка задачи}

Задача question answering (QA) состоит в автоматическом формировании ответа
на естественно-языковой вопрос на основе текстовых документов.

Рассматриваются два сценария. В режиме \emph{oracle} системе передаётся
гарантированно релевантный фрагмент, что позволяет оценить верхнюю границу
качества при идеальном ретривере. В режиме \emph{retriever-augmented} (RAG)
контекст формируется автоматически с помощью ранкеров (top-1 или top-5).

Таким образом, задача разбивается на две части: поиск релевантного документа
и генерация ответа.

\section{Описание данных}

Коллекция MIRAGE содержит 7560 пар вопрос--ответ из нескольких источников
(NaturalQA, PopQA, TriviaQA и др.). Каждому вопросу сопоставлен список
эталонных ответов, а также пул из 5 кандидатных документов, один из которых
является релевантным.

Пример запроса:
\begin{quote}
\textit{What is John Mayne's occupation?}
\end{quote}

Пример ответа:
\begin{quote}
\textit{journalist}
\end{quote}

\section{Методы и модели}

\subsection{Экстрактивный QA}

Используется модель \texttt{deepset/roberta-base-squad2}, обученная на SQuAD~2.0.
Модель предсказывает начало и конец ответа внутри документа.

В режиме top-5 модель применяется к каждому из фрагментов, после чего
выбирается ответ с максимальным confidence score.

\subsection{Генеративный QA}

Используется модель \texttt{Qwen2.5-1.5B-Instruct}. Ответ генерируется
по шаблону:

\begin{quote}
\texttt{Answer the question concisely based on the context.}
\end{quote}

Генерация выполняется жадным декодированием с ограничением длины ответа.

\subsection{Ранкеры}

Используются два ранкера:
\begin{itemize}
    \item BM25;
    \item CrossEncoder (\texttt{ms-marco-MiniLM-L-6-v2}).
\end{itemize}

Полнота (recall) релевантного документа:
\begin{itemize}
    \item BM25: 0.464 при top-1, 1.0 при top-5;
    \item CrossEncoder: 0.803 при top-1, 1.0 при top-5.
\end{itemize}

\subsection{Метрики}

Используются следующие метрики:
\begin{itemize}
    \item EM и F1 (SQuAD);
    \item EM-strict и EM-loose (MIRAGE);
    \item BERTScore.
\end{itemize}

\section{Результаты}

\begin{table*}[t]
\centering
\small
\begin{tabular}{llccccc}
\toprule
Модель & Контекст & EM & F1 & EM-strict & EM-loose & BERTScore \\
\midrule
RoBERTa (extractive) & Oracle      & 0.6392 & 0.7519 & 0.6392 & 0.8103 & 0.9522 \\
RoBERTa (extractive) & BM25 топ-1  & 0.3265 & 0.4076 & 0.3265 & 0.4459 & 0.9009 \\
RoBERTa (extractive) & CE топ-1    & 0.5396 & 0.6403 & 0.5396 & 0.6880 & 0.9392 \\
RoBERTa (extractive) & BM25 топ-5  & 0.5216 & 0.6119 & 0.5216 & 0.6512 & 0.9357 \\
RoBERTa (extractive) & CE топ-5    & 0.5216 & 0.6119 & 0.5216 & 0.6512 & 0.9357 \\
RoBERTa (extractive) & Mixed       & 0.5216 & 0.6119 & 0.5216 & 0.6512 & 0.9357 \\
\midrule
Qwen2.5-1.5B (gen)   & Oracle      & 0.4858 & 0.6608 & 0.4858 & 0.8634 & 0.9339 \\
Qwen2.5-1.5B (gen)   & BM25 топ-1  & 0.2464 & 0.3552 & 0.2464 & 0.4521 & 0.8921 \\
Qwen2.5-1.5B (gen)   & CE топ-1    & 0.4094 & 0.5629 & 0.4094 & 0.7287 & 0.9219 \\
Qwen2.5-1.5B (gen)   & BM25 топ-5  & 0.2061 & 0.3567 & 0.2061 & 0.6470 & 0.8906 \\
Qwen2.5-1.5B (gen)   & CE топ-5    & 0.2594 & 0.4291 & 0.2594 & 0.7646 & 0.9015 \\
Qwen2.5-1.5B (gen)   & Mixed       & 0.2844 & 0.4480 & 0.2844 & 0.7544 & 0.9043 \\
\bottomrule
\end{tabular}
\caption{Основные результаты экспериментов на MIRAGE}
\label{tab:main}
\end{table*}

\begin{table}[t]
\centering
\small
\begin{tabular}{lccccc}
\toprule
Конфигурация & LLM-judge & EM & EM-loose & F1 & BERTScore \\
\midrule
ext oracle   & 0.1312 & 0.6392 & 0.8103 & 0.7519 & 0.9522 \\
gen oracle   & 0.2546 & 0.4858 & 0.8634 & 0.6608 & 0.9339 \\
ext CE топ-1 & 0.1128 & 0.5396 & 0.6880 & 0.6403 & 0.9392 \\
gen CE топ-1 & 0.2232 & 0.4094 & 0.7287 & 0.5629 & 0.9219 \\
\bottomrule
\end{tabular}
\caption{LLM-as-a-judge vs автоматические метрики}
\label{tab:judge}
\end{table}

\section{Анализ результатов}

\subsection{Oracle vs retriever}

В oracle режиме достигается максимальное качество.
Экстрактивная модель показывает EM=0.6392, тогда как генеративная — 0.4858.
Разница объясняется тем, что извлечение точного ответа проще при наличии
релевантного документа.

\subsection{Влияние ранкера}

CrossEncoder значительно превосходит BM25 при top-1 (recall 0.803 против 0.464),
что приводит к заметному росту качества QA.

\subsection{Top-1 vs top-5}

Переход к top-5 повышает качество для экстрактивной модели,
так как релевантный документ почти всегда присутствует в наборе.

Для генеративной модели увеличение контекста не даёт прироста,
а иногда ухудшает результат из-за добавления нерелевантной информации.

\subsection{Экстрактивный vs генеративный подход}

Экстрактивная модель лучше по EM и F1.
Генеративная модель показывает более высокие значения EM-loose,
что связано с вариативностью формулировок ответов.

\subsection{LLM-as-a-judge}

Метрика LLM-as-a-judge слабо коррелирует с EM,
но ближе к EM-loose. Генеративная модель получает более высокие оценки,
что связано с семантической гибкостью ответов.

\section{Заключение}

В работе сравнивались экстрактивный и генеративный подходы к QA на MIRAGE.

Основные выводы:
\begin{itemize}
\item Экстрактивная модель обеспечивает наилучшее качество при наличии
      релевантного документа.
\item CrossEncoder значительно улучшает качество RAG по сравнению с BM25.
\item Использование top-5 полезно для экстрактивного QA, но не для генеративного.
\item Генеративные модели лучше отражают семантическое сходство,
      но уступают по точным метрикам.
\end{itemize}

Этот и другие отчeты, а также исходный код доступны в репозитории:
\url{https://github.com/qw1zzard/itmo-ir-search}.

\end{document}
